\section{Introduction}
\label{sec:intro}
Generative inpainting can add, remove, or replace image content while preserving much of the surrounding scene. Its accessibility makes detection and localization valuable forensic tasks: an analyst may need both an image-level decision and a map of the region supporting that decision. Our setting uses a single input photograph, without requiring the original image at inference.

Local edits are challenging because the modified region can occupy a small fraction of the image. Whole-image resizing provides scene context, but may suppress local details. Native-resolution crops retain detail, but each crop can lack context or contain no annotated edit. Natural textures, compression, and image processing can also produce strong residual responses. An image-level prediction alone does not establish that the detector identified the edited pixels; prior work shows that models can rely on cues outside the inpainted region~\cite{inpx}.

Manipulation localization methods have explored low-level forensic traces and pretrained visual representations~\cite{trufor,declip,opensdi}. FUSED combines forensic and semantic branches and reports cross-generator and transfer evaluations~\cite{fused}. Building on these directions, we study a detector with semantic and residual feature streams, whole-image and crop views, and separate outputs for image detection and localization. In a joint-view training configuration, a crop is labeled positive only if its annotation contains an edited pixel. The earlier ensemble used different checkpoints and training, and is evaluated separately.

Our contributions are: (1) a description of a single-image, multi-stream detector and its whole-image and tile-based inference; (2) an investigation of crop-aware supervision and the tradeoff between missed edits and false alarms; and (3) an error analysis by generator and edit size, with published baselines reported alongside the different evaluation conditions. The present development data and the 10K target scorecard do not establish a matched performance ranking against prior work.

\section{Related Work}
\label{sec:related}
\paragraph{Image manipulation localization.} Methods such as MVSS-Net, PSCC-Net, CAT-Net and TruFor investigate noise, compression, multi-scale features or learned residuals to identify edited regions~\cite{mvss,pscc,catnet,trufor}. These approaches demonstrate the value of local evidence, while genuine image processing can also create residual responses.

\paragraph{Generative image detection.} CNNDetection, GramNet, NPR and frequency-based approaches examine statistical clues associated with image generation~\cite{cnndet,gramnet,npr,freqnet}. Such image-level signals may generalize differently across generators and do not necessarily identify the modified region. The inpainting exchange study shows that detectors can rely on artifacts outside the edit~\cite{inpx}.

\paragraph{Pretrained representations and diffusion inpainting.} DeCLIP decodes CLIP features for localization, while OpenSDI studies diffusion-generated edits and introduces MaskCLIP~\cite{declip,opensdi}. FUSED combines forensic and semantic representations with token-level mixture-of-experts fusion and reports cross-generator results and transfer tests~\cite{fused}. Our work examines a different implementation using ConvNeXt, DINOv2, residual filtering and a cross-attention pyramid. Combining semantic and forensic features is an established direction; the contribution of any specific component in our configuration requires controlled ablation.

\section{Methodology}
\label{sec:method}
Given an RGB image $I\in\mathbb{R}^{H\times W\times3}$, the system predicts an image-level edit score and a spatial map. The original photograph is not needed at inference. Reference labels and masks are used during training and evaluation. Figure~\ref{fig:framework_flowchart} summarizes the documented ensemble evaluation; its compact and larger models use different backbones and fusion modules.
\begin{figure*}[t]
\centering
\includegraphics[width=\textwidth]{figures/framework_flowchart.pdf}
\caption{Documented development ensemble pipeline. A whole-image view supplies context and native tiles preserve local detail. The Base model uses cross-attention; the Small model uses gated fusion. The displayed decision settings are specific to the saved ensemble evaluation.}
\Description{A whole image and overlapping crops are processed through semantic and residual features, followed by mask prediction and score aggregation.}
\label{fig:framework_flowchart}
\end{figure*}

\subsection{Content and residual encoders}
The larger detector uses ConvNeXt-Base and frozen DINOv2 ViT-Base/14~\cite{convnext,dinov2}. ConvNeXt supplies features at strides 4, 8, 16 and 32. DINOv2 features from intermediate blocks 2, 5, 8 and 11 are interpolated to the corresponding spatial sizes. The two feature streams are concatenated and projected to 128 channels at each scale. The residual encoder combines 30 fixed SRM responses with 30 learned constrained-convolution responses~\cite{bayar,srm}, then produces features at strides 2, 4, 8 and 16.

At stride 16 the larger model uses eight-head cross-attention: projected forensic features provide queries and projected content features provide keys and values. The attention output includes a residual connection and feed-forward block. A learned spatial gate combines content and attended features; a top-down pyramid and further gates fuse the stride-8 and stride-4 representations. Stride-2 residual features refine the localization map, which is bilinearly upsampled. The compact model uses ConvNeXt-Small and frozen DINOv2-Small with gated fusion.

The image head concatenates global average and maximum pooled stride-32 content features with the mean of the highest 2,048 mask logits from that forward pass. An MLP converts these features to a classification logit. This internal logit pooling is distinct from the top-100 probability pooling used to aggregate the full tile map.

\subsection{Training configurations}
The compact and larger checkpoints used in the earlier ensemble were trained with an earlier crop-based procedure that could assign the source image label to a crop containing no labeled edit. A later, separate joint-view configuration gives the resized whole image its source label and assigns each native crop a positive label only if its mask contains edited pixels. The crop sampler favors edited regions for some training samples, while the mini-batch sampler balances genuine and edited source images.

The joint-view objective is
\begin{equation}
\mathcal{L}=\mathcal{L}_{\mathrm{whole}}+0.5\mathcal{L}_{\mathrm{crop}}+\mathcal{L}_{\mathrm{seg}}+0.3(\mathcal{L}_{\mathrm{bg,w}}+\mathcal{L}_{\mathrm{bg,c}}).
\end{equation}
The classification terms use binary cross-entropy. Segmentation on edit-containing crops combines balanced pixel binary cross-entropy and soft Dice loss; the background terms penalize mean mask probabilities on genuine whole views and empty-mask crops. Checkpoint selection in this experiment uses validation whole-image balanced accuracy. The saved ensemble results are not evidence of the joint-training objective's isolated effect.

\subsection{Development ensemble inference}
The recorded ensemble uses a $384\times384$ whole-image view and overlapping $384\times384$ tiles with stride 256. The larger and compact model predictions are combined with weights 0.65 and 0.35. Tile predictions include horizontal flips and Hann-weighted merging. The mean of the 100 highest values in the merged probability map gives $s_{\mathrm{local}}$, which is combined with the whole-image score:
\begin{equation}
 s = 0.15s_{\mathrm{whole}}+0.85s_{\mathrm{local}}.
\end{equation}
For the saved evaluation, the image threshold was 0.9358, inherited from an earlier calibration, and the mask threshold was 0.35. The mask is further processed by filling holes and removing components smaller than 64 pixels. These settings characterize the development evaluation; they have not been established for a separate claimed 10K run.

\section{Experiments and Results}
\label{sec:results}

\subsection{Data and evaluation}
The saved development evaluation uses 100 previously inspected images from OpenSDID-derived subsets~\cite{opensdi}: 20 images for each of SD1.5, SD2.1, SDXL, SD3 and FLUX.1, with ten genuine and ten edited in each group. The subset informed development and is not an independent final test. A prior audit found no exact hash overlap in the inspected manifests; this does not prove source disjointness or exclude near duplicates. Reference mask semantics and preprocessing should be checked before a definitive benchmark.

Image accuracy measures the fraction of correct decisions; edited-class F1 summarizes precision and recall, and AUROC assesses score ranking across thresholds. On edited images, localization IoU is the intersection of the predicted and reference masks divided by their union, averaged per image. Pixel F1 can instead be pooled over pixels or averaged over images; these definitions need not agree. Size bins use the fraction of the image covered by the reference mask: below 2\%, 2--10\%, and above 10\%.

\subsection{Recorded development evaluation}
A saved 100-image ensemble summary reports 70.0\% accuracy, 62.50\% image F1, 0.7452 AUROC, and 34.89\% mean IoU. It reports pooled pixel F1 of 58.30\% over edited images. The recorded ensemble uses 384-pixel views, tile stride 256, image score weights 0.15/0.85, a mask threshold of 0.35 and an inherited image threshold of 0.9358. These settings are specific to that evaluation. Table~\ref{tab:dev} shows the recorded per-generator results. Each accuracy uses 20 images and each IoU uses ten edited images.
\begin{table}[t]
\centering\small
\caption{Saved ensemble development results by editor. Accuracy and IoU are percentages.}
\label{tab:dev}
\begin{tabular}{@{}lrrr@{}}\toprule
Editor & Accuracy & AUROC & Mean IoU \\\midrule
SD1.5 & 75.0 & 0.9000 & 55.03 \\
SD2.1 & 70.0 & 0.7100 & 35.84 \\
SDXL & 60.0 & 0.5600 & 24.57 \\
SD3 & 70.0 & 0.8700 & 33.43 \\
FLUX.1 & 75.0 & 0.7100 & 25.58 \\\bottomrule
\end{tabular}
\end{table}
The reported IoU for tiny, medium and large edits was 8.94\%, 29.33\%, and 45.16\%, respectively (Figure~\ref{fig:edit-size}). These exploratory values indicate a need to inspect small-edit failures; their small bin sizes limit precision. A separate whole-image comparison recorded 13 false positives and 23 false negatives for the earlier Base model, versus eight and 33 for the joint-view model, each on 50 genuine and 50 edited images at its own validation-selected threshold. Different thresholds and training configurations prevent attributing the difference to any one component.
\begin{figure}[t]
\centering
\includegraphics[width=\columnwidth]{figures/edit-size-iou.png}
\caption{Saved ensemble development IoU by reference edit area.}
\Description{Three bars show the development mean IoU for tiny, medium, and large edits.}
\label{fig:edit-size}
\end{figure}

\subsection{10K target scorecard and benchmark context}
A local planning log labels the following numbers a target scorecard for a proposed 10,000-image run with 15,000 steps, a global batch of 16, and two Tesla T4 GPUs. Table~\ref{tab:reported10k} preserves its aggregate figures beside published FUSED results~\cite{fused}. The 10K checkpoint, training manifest, prediction file, and full aggregation rules have not been established from the available artifacts. FUSED reports equal-weight averages across five generator groups after SD1.5-only training, while the 10K protocol remains incompletely documented; the columns should not be interpreted as a matched test.
\begin{table}[t]
\centering\small
\caption{Published FUSED OpenSDID averages and the separately reported 10K target. Protocols and aggregation are not established as matched.}
\label{tab:reported10k}
\begin{tabular}{@{}lrr@{}}\toprule
Metric (\%) & FUSED & 10K target \\\midrule
Image accuracy & 89.50 & 91.20 \\
Image F1 & 88.00 & 89.60 \\
Mean localization IoU & 44.50 & 47.30 \\
Pixel F1 & 55.30 & 63.80 \\\bottomrule
\end{tabular}
\end{table}
FUSED uses per-image pixel F1; the aggregation of the reported 10K target pixel F1 is unknown. The 10K target records 46 true positives, 47 true negatives, three false positives and four false negatives. For those 100 predictions, the counts imply 93.0\% accuracy, 93.0\% balanced accuracy and 92.93\% edited-class F1, rather than the reported 91.20\%, 92.60\% and 89.60\%. The five reported per-generator IoUs have an unweighted mean of 49.70\%, rather than the reported 47.30\%. Different subsets or weighting might explain the mismatch, but cannot be established without the underlying predictions and evaluation specification.

\subsection{Published baseline context}
Table~\ref{tab:baselines} lists the methods and reported scores in the common OpenSDID evaluation of FUSED~\cite{fused}. Accuracy is image-level; IoU is localization overlap. The venue refers to each method's original publication, while the measurements come from FUSED. These literature scores were not obtained in the 100-image development evaluation described above.
\begin{table}[t]
\centering\small
\caption{Published OpenSDID averages, in percent, as reported by FUSED~\cite{fused}. Dashes denote unavailable localization scores.}
\label{tab:baselines}
\begin{tabular}{@{}llrr@{}}\toprule
Method & Venue & Acc. & IoU \\\midrule
FUSED~\cite{fused} & arXiv 2026 & 89.5 & 44.5 \\
MaskCLIP~\cite{opensdi} & CVPR 2025 & 82.0 & 42.7 \\
RINE~\cite{rine} & ECCV 2024 & 80.3 & -- \\
PSCC-Net~\cite{psccnet} & TCSVT 2022 & 76.8 & 29.7 \\
NPR~\cite{npr} & CVPR 2024 & 76.4 & -- \\
CAT-Net~\cite{catnet} & IJCV 2022 & 76.2 & 37.4 \\
UniFD~\cite{unifd} & CVPR 2023 & 75.7 & -- \\
MVSS-Net~\cite{mvss} & TPAMI 2023 & 75.1 & 29.8 \\
DeCLIP~\cite{declip} & WACV 2025 & 73.1 & 25.2 \\
GramNet~\cite{gramnet} & CVPR 2020 & 72.3 & -- \\
CNNDetection~\cite{cnnspot} & CVPR 2020 & 70.9 & -- \\
TruFor~\cite{trufor} & CVPR 2023 & 69.8 & 36.9 \\
ObjectFormer~\cite{objectformer} & CVPR 2022 & 66.3 & 24.1 \\
FreqNet~\cite{freqnet} & AAAI 2024 & 66.3 & -- \\
IML-ViT~\cite{imlvit} & arXiv 2023 & 56.4 & 32.5 \\\bottomrule
\end{tabular}
\end{table}

\subsection{Qualitative examples}
Figure~\ref{fig:qual} presents one earlier diagnostic prediction, showing partial overlap with the reference edit alongside responses on genuine background regions. It is a selected example from an earlier configuration, not a prediction attributed to the 10K target or the development ensemble.
\begin{figure*}[t]
\centering
\includegraphics[width=\textwidth]{figures/sd15-localization.png}
\caption{Earlier SD1.5 diagnostic example: input, reference region, anomaly heatmap, and prediction overlay. Yellow indicates overlap with the reference; red denotes predictions outside it.}
\Description{An image, reference edit, anomaly heatmap, and prediction overlay illustrate partial localization and background responses.}
\label{fig:qual}
\end{figure*}


\section{Discussion and Limitations}
Our results show that detection and localization must be considered together: an image can receive the correct label while its predicted mask misses much of the edit. The low overlap for small regions and the differing false-alarm and missed-edit counts motivate evaluating crop supervision, tile inference, and decision thresholds in controlled comparisons. Existing configurations differ in training and scoring, so their differences cannot isolate the contribution of a single component.

The development subset was inspected repeatedly and is small. Hash-based checks alone do not establish source-level independence, and mask provenance and transformations require verification. A stronger evaluation should lock the source-group split and thresholds before testing, report uncertainty and edit-size sample counts, and evaluate common checkpoints on the same images. Further tests should assess compression, resizing, held-out editors, and paired controls for reliance on traces outside the edit~\cite{inpx}. The claimed 10K run requires its checkpoint, manifest, full logs, predictions, and aggregation rules before it can support a matched comparison. Hardware latency and component ablations also require measured runs under documented settings.

\section{Conclusion}
We studied a dual-view approach to single-image inpainting detection and localization. The development evaluation suggests that small edited regions and the balance between missed edits and false alarms remain important challenges. The 10K scorecard is retained as a target, not a verified experimental result. A source-disjoint evaluation with reproducible predictions and matched baselines is required to determine whether the approach improves on published methods.
